\section{Mixed-radix hierarchy budgets and a conditional complexity contract}
\label{sec:hierarchy}

A hierarchy algorithm can terminate by repeatedly simplifying an early surface
and rebuilding the later surfaces.  Lexicographic descent proves finiteness,
but the quantitative question is how many such rebuilds occur and how much
work each requires.  We refine the abstract accounting to unequal per-level
budgets and give an exact formula for suffix reconstruction.  These are
combinatorial results about a stated transition model.  An implementation must
separately establish that its geometric operations satisfy that model.

\subsection{Uniform per-level caps}

\begin{definition}[A bounded hierarchy epoch]
An epoch has a fixed maximum length $L$ and integer caps $g_i\geq0$,
$1\leq i\leq L$.  Whenever its construction phase returns, it records a
completed, zero-padded vector $x=(x_1,\ldots,x_L)$ with $0\leq x_i\leq g_i$.
Every nonterminal rollback produces a strictly smaller vector at the next
return, in lexicographic order with the first coordinate most significant.
The caps hold uniformly over every allowed suffix reconstruction in the
epoch.
\end{definition}

Let $r_i=g_i+1$, let $P=\prod_{i=1}^L r_i$, and define the mixed-radix rank
\begin{equation}
 \Phi(x)=\sum_{i=1}^L x_i\prod_{h=i+1}^L r_h.
 \label{eq:hierarchy-rank}
\end{equation}
Empty products equal one.  A radix equal to one represents a coordinate fixed
at zero; it is allowed throughout the exact counting statements.

\begin{proposition}[Mixed-radix descent]
\label{prop:mixed-radix}
Within a bounded hierarchy epoch there are at most $P-1$ strictly decreasing
transitions between completed vectors.  This bound is sharp for the abstract
model.
\end{proposition}

\begin{proof}
The rank lies between zero and $P-1$.  If coordinate $j$ is the first to
decrease, its contribution drops by at least $\prod_{h>j}r_h$.  The maximum
possible increase of all later coordinates is
\[
 \sum_{i>j}(r_i-1)\prod_{h>i}r_h
   =\prod_{h>j}r_h-1.
\]
Earlier contributions are unchanged, so $\Phi$ drops by at least one.  The
ordinary mixed-radix countdown from the largest vector to zero attains all
$P$ ranks and has $P-1$ transitions.
\end{proof}

For equal caps this recovers the familiar product $(g+1)^L$ appearing in the
iteration analysis of Lackenby's hierarchy algorithm~\cite{lackenby2026}.
Unequal caps matter because an early level may permit many simplifications
while a later level permits few.  Replacing every cap by the largest one can
discard useful information.  Nevertheless, an observed maximum from one
successful run is not a uniform cap: rebuilding an earlier prefix can expose
larger complexity at a later level.

A guaranteed decrement can also improve the ranking.  Suppose coordinate $i$
drops by at least an integer $d_i\geq1$ whenever it is the first changed
coordinate.  Replacing it by $\lfloor x_i/d_i\rfloor$ gives a strict decrease
with radix $1+\lfloor g_i/d_i\rfloor$.  This replacement requires a decrement
guarantee for every relevant transition, including those following rebuilt
suffixes; a typical decrement or an empirical average is insufficient.

\subsection{The exact cost of rebuilding affected suffixes}

The crude estimate charges $L$ visits after every decrease.  If a rollback
first changing coordinate $j$ only requires revisiting levels $j,\ldots,L$,
one can count the charged visits more accurately.  This remains meaningful
even when some coordinates have radix one, because a fixed coordinate may
still represent a physical construction step.

\begin{theorem}[Sharp suffix reconstruction budget]
\label{thm:suffix-budget}
Let $r_1,\ldots,r_L$ be arbitrary positive integer radices and set
\[
 R_0=1,\qquad R_i=\prod_{h=1}^i r_h,\qquad P=R_L.
\]
Give initial construction charge at most $L$.  Give a strict lexicographic
decrease whose first changed coordinate is $j$ charge at most $L-j+1$.
Every decreasing trace then has total charge at most
\begin{equation}
 W=\sum_{i=1}^L R_i
  =L+\sum_{j=1}^L(L-j+1)(r_j-1)R_{j-1}.
 \label{eq:exact-suffix-budget}
\end{equation}
The full mixed-radix countdown with all allowed charges attains equality.
If every $r_i\geq2$, then $W<2P$.  Without that additional condition,
$W\leq LP$ always holds.
\end{theorem}

\begin{proof}
Consider a jump from vector $x$ to a smaller vector $y$, first changing
coordinate $j$.  The consecutive countdown transitions between their ranks
include a transition first changing that same coordinate: their earlier
coordinates agree and their $j$th coordinates differ.  Select one such
transition.  Jumps in a strictly decreasing trace occupy disjoint rank
intervals, so the selected countdown transitions are distinct.  Each selected
transition has charge $L-j+1$, at least the charge assigned to its jump.
Therefore the trace is dominated by the full countdown, including the same
initial construction allowance.

For each assignment of the first $j-1$ digits, the full countdown decreases
digit $j$ exactly $r_j-1$ times while keeping those earlier digits fixed.
There are $R_{j-1}$ such assignments.  Hence the number of transitions whose
first changed digit is $j$ is
$(r_j-1)R_{j-1}=R_j-R_{j-1}$.  This proves the second expression in
\eqref{eq:exact-suffix-budget}.  Telescoping gives the first expression:
\[
 L+\sum_{j=1}^L(L-j+1)(R_j-R_{j-1})=\sum_{i=1}^L R_i.
\]
When all radices are at least two, $R_i\leq P/2^{L-i}$ and a finite
geometric sum is strictly less than $2P$.  In general $R_i\leq P$, giving
$W\leq LP$.  The full countdown realizes the counted transitions, proving
sharpness.
\end{proof}

For example, radices $(3,1,2)$ give $P=6$ and $W=3+3+6=12$.
The exact budget equals $2P$, so a strict $2P$ estimate would be false here.
If all radices are one, there is a single state but initial construction still
costs $L$.  Removing fixed digits from the rank does not automatically remove
their work from the algorithm.

\begin{corollary}[Unequal costs per level]
\label{cor:weighted-suffix}
Let level $i$ have a nonnegative charge allowance $a_i$.  If initialization
costs at most $\sum_i a_i$ and a jump first changing $j$ costs at most
$\sum_{i=j}^L a_i$, then the total charge is at most
\begin{equation}
 \sum_{i=1}^L a_iR_i.
 \label{eq:weighted-suffix}
\end{equation}
This is also sharp in the abstract model.
\end{corollary}

\begin{proof}
Use the same disjoint-interval injection into the countdown.  Its total is
\[
 \sum_{i=1}^L a_i+
 \sum_{j=1}^L(R_j-R_{j-1})\sum_{i=j}^L a_i
 =\sum_{i=1}^L a_iR_i.
\]
The countdown with every allowance charged attains equality.
\end{proof}

This weighted form isolates another optimization target: reducing the cost of
late levels can matter more because they are reconstructed more frequently.
The statement applies only when the supplied allowances bound the entire
work assigned to a visit, rather than the cost of one selected subroutine.

\subsection{A conditional theorem with explicit implementation obligations}

\begin{theorem}[Conditional recognition complexity]
\label{thm:conditional-qp}
Suppose a sound and complete recognizer on inputs of size $n$ satisfies the
following conditions:
\begin{enumerate}[label=(\roman*)]
\item It has at most $E(n)$ epochs, each with length at most $L(n)$ and
uniform caps bounded by $g_i(n)$, and satisfies
Proposition~\ref{prop:mixed-radix} in each epoch.
\item All initialization and rollback work can be assigned to charged level
visits as in Theorem~\ref{thm:suffix-budget}.  Each visit has total bit-cost
at most $C(n)$, including repeated calls, surface searches, normalization,
geometric updates, and necessary bookkeeping.
\item Work outside those visits has total cost at most $T_0(n)$.  All stored
geometric data and metadata have a specified bit-size bound $B(n)$ under
which the preceding operation bounds hold.
\end{enumerate}
Put $r_i(n)=g_i(n)+1$ and $P(n)=\prod_i r_i(n)$.  Then
\begin{equation}
 T(n)\leq T_0(n)+O\!\left(E(n)C(n)\sum_{i=1}^{L(n)}
                       \prod_{h=1}^i r_h(n)\right).
 \label{eq:conditional-runtime}
\end{equation}
In particular this is at most $T_0+O(ECLP)$, and at most
$T_0+O(ECP)$ if all radices are at least two.  If $E,C,L,B,T_0$ are
polynomially bounded and
\begin{equation}
 H(n):=\sum_{i=1}^{L(n)}\log(g_i(n)+1)=O(\log^2 n),
 \label{eq:entropy-contract}
\end{equation}
then the recognizer runs in $n^{O(\log n)}$ time.
\end{theorem}

\begin{proof}
Apply Theorem~\ref{thm:suffix-budget} to each epoch and multiply the visit
count by the total cost allowance.  Sum over epochs and add $T_0$.  Finally,
$P=\exp H$; polynomial prefactors contribute only $O(\log n)$ to the
logarithm of the runtime.
\end{proof}

If only Proposition~\ref{prop:mixed-radix} is known, the conservative $LP$
estimate still follows provided each rollback and initial construction has
total cost at most $LC$.  Either formulation must control complete work,
not merely one primitive's runtime.  A polynomial-time routine called an
unbounded number of times does not satisfy the theorem by itself.

The sum $H$, rather than maximum hierarchy depth alone, is the relevant
budget.  A logarithmic number of levels with polynomial caps gives
$H=O(\log^2 n)$.  A squared-logarithmic number of levels with polynomial
caps gives only $H=O(\log^3 n)$ and therefore
$2^{O(\log^3 n)}$ time.  Squared-logarithmic depth with uniformly bounded
caps would again give the sharper $n^{O(\log n)}$ target.  Both displayed
time classes are quasi-polynomial; the exponent of the logarithm distinguishes
the stronger target.  Removing a multiplicative $L$ through amortized suffix
accounting does not itself change that exponent when $L$ is polynomial.

\subsection{Relation to the published hierarchy description}

Lackenby's February 2021 slides announce $n^{O(\log n)}$
time~\cite{lackenbyfeb}.  The March Oxford slides instead state
$2^{O(\log^3 n)}$, with a squared-logarithmic multisurface-depth
bound~\cite{lackenbymarch}.  These should be recorded as distinct quantitative
announcements.

The July 2026 preprint gives a full hierarchy algorithm but does not state a
quasi-polynomial runtime theorem.  Its Proposition~9.1 gives the
$L(g+1)^L$ iteration bound; Proposition~10.2 bounds the refined ordinary
hierarchy length by $4c_q(H)$; Theorem~14.1 gives polynomially encoded and
verifiable essential hierarchies.  Section~9 explicitly leaves some primitive
running times unspecified~\cite{lackenby2026}.  A short certificate and a
verification procedure do not bound deterministic search for that certificate.

At source level, Step~15 discards the surfaces following the one being
simplified.  Steps~12--14 also transport and modify a violating disc while
moving through the hierarchy~\cite{lackenby2026}.  Suffix deletion therefore
motivates our accounting model, but is not a proof that every geometric cost
fits it.  Theorem~\ref{thm:conditional-qp} deliberately retains that
obligation.  Additional outer restarts would likewise need their own epoch
bound, rather than inheriting one from the inner rank.

Compressed representations are central to the bit-cost condition.  A normal
surface can have many disks while its coordinates have short binary encodings.
The Agol--Hass--Thurston orbit-counting method supplies polynomial algorithms
in the logarithm of represented interval sizes, including compressed normal
surface component counting~\cite{aht}.  Expanding every represented disk or
parallel component would destroy that advantage.  A complete implementation
must specify the representations passed between primitives and show that
cutting, normalization, and disc transport preserve their certified size
bounds.

\subsection{An independently executable finite audit}

The program \path{tools/hierarchy_budget.py} contains no geometric backend.
For a small radix tuple it constructs the directed acyclic graph containing
every strict lexicographic transition, assigns each edge its affected-suffix
charge, and computes a maximum-weight path.  This calculation does not assume
that consecutive countdown transitions are optimal.  It checks that the
computed maximum, the explicit full-countdown charge, and
$\sum_iR_i$ agree.  For very small state spaces it also enumerates every
nonempty decreasing sequence directly.

The supplied audit covers all 364 radix tuples of length at most five over
$\{1,2,3\}$.  It checks 284,932 possible strict transitions and explicitly
enumerates 17,021 decreasing sequences in its smallest cases; all agree with
the theorem.  Results appear in \path{data/hierarchy_budget_audit.json}.
These finite checks support reproducibility and catch indexing or unit-radix
mistakes.  The proof above establishes the unrestricted combinatorial result;
neither the proof nor the program supplies the missing topological charging
or search bounds.
